% ==============================================================================
% PRESENTACIÓN LaTeX (Beamer) — Seguridad de BD ArtiSync
% Universidad Técnica Estatal de Quevedo
% ==============================================================================
\documentclass[aspectratio=169]{beamer}

% --- Tema y colores ---
\usetheme{Madrid}
\usecolortheme{whale}

% Colores personalizados UTEQ
\definecolor{uteqgreen}{RGB}{0,100,60}
\definecolor{uteqdark}{RGB}{20,40,30}
\definecolor{uteqlight}{RGB}{220,245,230}
\definecolor{uteqaccent}{RGB}{0,150,90}
\definecolor{sqlbg}{RGB}{245,245,250}
\definecolor{sqlkey}{RGB}{0,0,170}
\definecolor{sqlcom}{RGB}{0,128,0}

\setbeamercolor{palette primary}{bg=uteqgreen,fg=white}
\setbeamercolor{palette secondary}{bg=uteqdark,fg=white}
\setbeamercolor{palette tertiary}{bg=uteqgreen,fg=white}
\setbeamercolor{palette quaternary}{bg=uteqdark,fg=white}
\setbeamercolor{structure}{fg=uteqgreen}
\setbeamercolor{title}{fg=white,bg=uteqgreen}
\setbeamercolor{frametitle}{fg=white,bg=uteqgreen}
\setbeamercolor{block title}{bg=uteqgreen,fg=white}
\setbeamercolor{block body}{bg=uteqlight,fg=black}
\setbeamercolor{block title alerted}{bg=red!70!black,fg=white}
\setbeamercolor{block body alerted}{bg=red!10,fg=black}
\setbeamercolor{block title example}{bg=uteqaccent,fg=white}
\setbeamercolor{block body example}{bg=green!5,fg=black}
\setbeamercolor{item}{fg=uteqgreen}

% --- Paquetes ---
\usepackage[utf8]{inputenc}
\usepackage[T1]{fontenc}
\usepackage[spanish]{babel}
\usepackage{amsmath}
\usepackage{graphicx}
\usepackage{booktabs}
\usepackage{array}
\usepackage{listings}
\usepackage{fontawesome5}
\usepackage{tikz}
\usetikzlibrary{shapes.geometric, arrows.meta, positioning, fit, calc}
\usepackage{multicol}

% --- Configuración SQL ---
\lstdefinestyle{sqlbeamer}{
    language=SQL,
    backgroundcolor=\color{sqlbg},
    basicstyle=\ttfamily\tiny,
    keywordstyle=\color{sqlkey}\bfseries,
    commentstyle=\color{sqlcom}\itshape,
    stringstyle=\color{red!70!black},
    frame=single,
    rulecolor=\color{gray!40},
    breaklines=true,
    showstringspaces=false,
    tabsize=2,
    morekeywords={SERIAL,BOOLEAN,TEXT,TIMESTAMP,DECIMAL,VARCHAR,
                  REFERENCES,CASCADE,UNIQUE,DEFAULT,NOLOGIN,LOGIN,
                  NOSUPERUSER,NOCREATEDB,NOCREATEROLE,NOINHERIT,
                  CREATEDB,CREATEROLE,INHERIT,GRANT,REVOKE,ROLE,
                  PASSWORD,CONNECTION,LIMIT,VALID,UNTIL,USAGE,
                  SCHEMA,ALL,PRIVILEGES,ON,TO,CONNECT,OWNER,
                  ENCODING,TABLESPACE,COMMENT}
}
\lstset{style=sqlbeamer}

% --- Quitar navegación ---
\setbeamertemplate{navigation symbols}{}

% --- Pie de página ---
\setbeamertemplate{footline}{
    \leavevmode%
    \hbox{%
        \begin{beamercolorbox}[wd=.333333\paperwidth,ht=2.25ex,dp=1ex,center]{palette primary}%
            \usebeamerfont{author in head/foot}\insertshortauthor
        \end{beamercolorbox}%
        \begin{beamercolorbox}[wd=.333333\paperwidth,ht=2.25ex,dp=1ex,center]{palette secondary}%
            \usebeamerfont{title in head/foot}\insertshorttitle
        \end{beamercolorbox}%
        \begin{beamercolorbox}[wd=.333333\paperwidth,ht=2.25ex,dp=1ex,right]{palette tertiary}%
            \usebeamerfont{date in head/foot}\insertshortdate{}\hspace*{2em}
            \insertframenumber{} / \inserttotalframenumber\hspace*{2ex}
        \end{beamercolorbox}}%
    \vskip0pt%
}

% --- Datos ---
\title[Seguridad BD ArtiSync]{Análisis de Seguridad de Base de Datos}
\subtitle{Usuarios, Roles y Privilegios en PostgreSQL --- ArtiSync}
\author[Bone, Carvajal, Figueroa, Rios]{
    Bone Niurca \and Carvajal Johan \and Figueroa Bryan \and Rios Jhon
}
\institute[UTEQ]{
    Facultad de Ciencias de la Ingeniería\\
    Universidad Técnica Estatal de Quevedo
}
\date{\today}

% ==============================================================================
\begin{document}

% --- PORTADA ---
\begin{frame}
\titlepage
\end{frame}

% --- ÍNDICE ---
\begin{frame}{Contenido}
\tableofcontents
\end{frame}

% ==============================================================================
% SECCIÓN 1: INTRODUCCIÓN
% ==============================================================================
\section{Introducción}

\begin{frame}{¿Qué es ArtiSync?}
    \begin{columns}[T]
        \begin{column}{0.55\textwidth}
            \begin{block}{Plataforma de Marketplace Creativo}
                ArtiSync conecta \textbf{artistas digitales} con \textbf{clientes} que buscan servicios creativos personalizados.
            \end{block}
            \vspace{0.3cm}
            \textbf{Stack tecnológico:}
            \begin{itemize}
                \item \faJava\ Java 21 + Spring Boot 3.2
                \item \faAngular\ Angular 17+
                \item \faDatabase\ PostgreSQL 16
                \item \faDocker\ Docker Compose
                \item \faKey\ JWT + Spring Security 6
            \end{itemize}
        \end{column}
        \begin{column}{0.42\textwidth}
            \begin{exampleblock}{7 Módulos Funcionales}
                \scriptsize
                \begin{enumerate}
                    \item Seguridad y Control de Acceso
                    \item Perfiles y Portafolio
                    \item Catálogo de Servicios
                    \item Pedidos y Flujos de Trabajo
                    \item Legal y Finanzas (Escrow)
                    \item Comunicación
                    \item Social y Comunidad
                \end{enumerate}
            \end{exampleblock}
        \end{column}
    \end{columns}
\end{frame}

\begin{frame}{Objetivos del Análisis}
    \begin{block}{Objetivo General}
        Diseñar e implementar el modelo de seguridad a nivel de base de datos para ArtiSync aplicando el \textbf{principio de mínimo privilegio}.
    \end{block}
    \vspace{0.3cm}
    \begin{columns}[T]
        \begin{column}{0.48\textwidth}
            \begin{enumerate}
                \item[\faUsers] Identificar los usuarios del sistema.
                \item[\faUserShield] Definir roles para cada usuario.
            \end{enumerate}
        \end{column}
        \begin{column}{0.48\textwidth}
            \begin{enumerate}
                \item[\faCog] Establecer privilegios del sistema.
                \item[\faLock] Definir privilegios sobre objetos.
            \end{enumerate}
        \end{column}
    \end{columns}
    \vspace{0.3cm}
    \begin{alertblock}{Principio Rector}
        \centering \textbf{Principio de Mínimo Privilegio (PoLP):} cada rol recibe \textit{únicamente} los permisos necesarios para cumplir su función.
    \end{alertblock}
\end{frame}

% ==============================================================================
% SECCIÓN 2: BASE DE DATOS
% ==============================================================================
\section{Estructura de la Base de Datos}

\begin{frame}{Base de Datos: 37 Tablas en 7 Módulos}
    \centering
    \scriptsize
    \begin{tabular}{>{\raggedright\arraybackslash}p{3.2cm} >{\raggedright\arraybackslash}p{9cm} c}
        \toprule
        \textbf{Módulo} & \textbf{Tablas principales} & \textbf{Cant.} \\
        \midrule
        \faShieldAlt\ Seguridad & \texttt{roles, permisos, usuarios, sesiones\_usuario, tokens\_recuperacion, autenticacion\_dos\_factores} & 9 \\
        \midrule
        \faUserCircle\ Perfiles & \texttt{perfiles\_creadores, certificados\_ia, habilidades, portafolios, portafolio\_items} & 7 \\
        \midrule
        \faStore\ Catálogo & \texttt{categorias, subcategorias, servicios, atributos\_dinamicos, etiquetas} & 7 \\
        \midrule
        \faClipboardList\ Pedidos & \texttt{flujos\_trabajo, etapas\_flujo, pedidos, historial\_estados\_pedido, tickets\_revision} & 7 \\
        \midrule
        \faFileContract\ Finanzas & \texttt{contratos, pagos\_garantia, transacciones\_pago, entregables\_finales} & 5 \\
        \midrule
        \faComments\ Comunicación & \texttt{salas\_chat, mensajes, documentos\_adjuntos, notificaciones\_sistema} & 5 \\
        \midrule
        \faHeart\ Social & \texttt{seguidores, comentarios\_portafolio, likes\_portafolio, resenas\_servicios, sorteos} & 6 \\
        \bottomrule
    \end{tabular}
\end{frame}

\begin{frame}{Creación de la Base de Datos}
    \begin{lstlisting}[title={Fase 0: CREATE DATABASE}]
CREATE DATABASE artisyncbd
    WITH
    OWNER = postgres
    ENCODING = 'UTF8'
    LC_COLLATE = 'es_EC.UTF-8'
    LC_CTYPE = 'es_EC.UTF-8'
    LOCALE_PROVIDER = 'libc'
    TABLESPACE = pg_default
    CONNECTION LIMIT = -1
    IS_TEMPLATE = False;

COMMENT ON DATABASE artisyncbd IS
    'Base de datos principal de la plataforma ArtiSync - '
    'Marketplace creativo para artistas digitales.';
    \end{lstlisting}
\end{frame}

% ==============================================================================
% SECCIÓN 3: USUARIOS
% ==============================================================================
\section{Identificación de Usuarios}

\begin{frame}{Usuarios del Sistema}
    \begin{columns}[T]
        \begin{column}{0.48\textwidth}
            \begin{block}{\faUsers\ Usuarios de Negocio (6)}
                \scriptsize
                \begin{description}
                    \item[\faUserCog\ Administrador] Superusuario. Gestiona usuarios, roles, catálogos y configuración global.
                    \item[\faPaintBrush\ Creador] Artista que ofrece servicios, gestiona perfil y portafolio.
                    \item[\faShoppingCart\ Cliente] Contrata servicios, realiza pedidos y evalúa trabajos.
                    \item[\faEye\ Moderador] Revisa certificados IA, modera comentarios y reportes.
                    \item[\faHeadset\ Soporte] Atiende incidencias consultando cuentas y pedidos.
                    \item[\faChartLine\ Auditor Fin.] Supervisa transacciones y genera reportes contables.
                \end{description}
            \end{block}
        \end{column}
        \begin{column}{0.48\textwidth}
            \begin{block}{\faRobot\ Usuarios Técnicos (3)}
                \scriptsize
                \begin{description}
                    \item[\faServer\ App Backend] Usuario no humano del servicio Spring Boot para operaciones CRUD vía JPA.
                    \item[\faHdd\ Backup] Tareas automatizadas de respaldo con \texttt{pg\_dump}.
                    \item[\faChartBar\ Solo Lectura] Herramientas de monitoreo y dashboards (Grafana, Metabase).
                \end{description}
            \end{block}
            \vspace{0.2cm}
            \begin{alertblock}{\faExclamationTriangle\ Total}
                \centering \textbf{9 tipos de usuarios} identificados
            \end{alertblock}
        \end{column}
    \end{columns}
\end{frame}

% ==============================================================================
% SECCIÓN 4: ROLES
% ==============================================================================
\section{Definición de Roles}

\begin{frame}{Modelo RBAC: Roles de Grupo y Usuarios de Login}
    \begin{columns}[T]
        \begin{column}{0.48\textwidth}
            \begin{block}{\faLayerGroup\ Roles de Grupo (\texttt{NOLOGIN})}
                \scriptsize
                Agrupaciones lógicas de permisos:
                \begin{enumerate}
                    \item \texttt{rol\_app\_backend}
                    \item \texttt{rol\_administrador}
                    \item \texttt{rol\_moderador}
                    \item \texttt{rol\_soporte}
                    \item \texttt{rol\_auditor\_fin}
                    \item \texttt{rol\_backup}
                    \item \texttt{rol\_solo\_lectura}
                \end{enumerate}
            \end{block}
        \end{column}
        \begin{column}{0.48\textwidth}
            \begin{block}{\faSignInAlt\ Usuarios de Login}
                \scriptsize
                Cuentas con conexión a la BD:
                \begin{enumerate}
                    \item \texttt{artisync\_app} $\rightarrow$ \texttt{rol\_app\_backend}
                    \item \texttt{artisync\_admin} $\rightarrow$ \texttt{rol\_administrador}
                    \item \texttt{artisync\_moderador} $\rightarrow$ \texttt{rol\_moderador}
                    \item \texttt{artisync\_soporte} $\rightarrow$ \texttt{rol\_soporte}
                    \item \texttt{artisync\_auditor} $\rightarrow$ \texttt{rol\_auditor\_fin}
                    \item \texttt{artisync\_backup} $\rightarrow$ \texttt{rol\_backup}
                    \item \texttt{artisync\_readonly} $\rightarrow$ \texttt{rol\_solo\_lectura}
                \end{enumerate}
            \end{block}
        \end{column}
    \end{columns}
\end{frame}

\begin{frame}{Diagrama de Jerarquía de Roles}
    \centering
    \begin{tikzpicture}[
        role/.style={rectangle, draw=uteqgreen!80, fill=uteqgreen!15, thick,
            minimum width=2.6cm, minimum height=0.65cm,
            font=\scriptsize\ttfamily, align=center, rounded corners=3pt},
        user/.style={rectangle, draw=blue!60, fill=blue!10, thick,
            minimum width=2.2cm, minimum height=0.55cm,
            font=\scriptsize\ttfamily, align=center, rounded corners=3pt},
        arrow/.style={-{Stealth[length=4pt]}, thick, draw=gray!60},
        inherit/.style={-{Stealth[length=5pt]}, very thick, draw=uteqaccent, dashed},
        node distance=0.5cm and 0.6cm
    ]
        % Roles de grupo - fila superior
        \node[role] (admin) {rol\_administrador};
        \node[role, below left=1.0cm and 0.6cm of admin] (app) {rol\_app\_backend};
        \node[role, below=1.0cm of admin] (mod) {rol\_moderador};
        \node[role, below right=1.0cm and 0.6cm of admin] (sop) {rol\_soporte};

        % Roles de grupo - fila inferior
        \node[role, below=1.0cm of app] (aud) {rol\_auditor\_fin};
        \node[role, below=1.0cm of mod] (bkp) {rol\_backup};
        \node[role, below=1.0cm of sop] (ro) {rol\_solo\_lectura};

        % Usuarios de login
        \node[user, below=0.6cm of admin] (ua) {\footnotesize artisync\_admin};
        \node[user, below=1.5cm of app] (uapp) {\footnotesize artisync\_app};
        \node[user, below=0.6cm of mod] (umod) {\footnotesize artisync\_moderador};
        \node[user, below=0.6cm of sop] (usop) {\footnotesize artisync\_soporte};
        \node[user, below=0.6cm of aud] (uaud) {\footnotesize artisync\_auditor};
        \node[user, below=0.6cm of bkp] (ubkp) {\footnotesize artisync\_backup};
        \node[user, below=0.6cm of ro] (uro) {\footnotesize artisync\_readonly};

        % Herencia
        \draw[inherit] (admin) -- node[left, font=\tiny\itshape, text=uteqaccent] {hereda} (app);

        % Usuarios -> Roles
        \draw[arrow] (ua) -- (admin);
        \draw[arrow] (uapp) -- (app);
        \draw[arrow] (umod) -- (mod);
        \draw[arrow] (usop) -- (sop);
        \draw[arrow] (uaud) -- (aud);
        \draw[arrow] (ubkp) -- (bkp);
        \draw[arrow] (uro) -- (ro);
    \end{tikzpicture}

    \vspace{0.2cm}
    \scriptsize
    \textcolor{uteqaccent}{\rule{1cm}{2pt}} Herencia de rol \quad
    \textcolor{gray}{\rule{1cm}{1pt}} Membresía de usuario
\end{frame}

\begin{frame}{Código: Creación de Roles de Grupo}
    \begin{lstlisting}[title={Roles de Grupo (NOLOGIN)}]
-- Rol para la aplicacion backend (Spring Boot / JPA)
CREATE ROLE rol_app_backend
    NOLOGIN NOSUPERUSER NOCREATEDB NOCREATEROLE NOINHERIT;

-- Rol de administrador (hereda del backend + DDL)
CREATE ROLE rol_administrador
    NOLOGIN NOSUPERUSER CREATEDB CREATEROLE INHERIT;
GRANT rol_app_backend TO rol_administrador;

-- Rol de moderador
CREATE ROLE rol_moderador
    NOLOGIN NOSUPERUSER NOCREATEDB NOCREATEROLE NOINHERIT;

-- Rol de soporte tecnico
CREATE ROLE rol_soporte
    NOLOGIN NOSUPERUSER NOCREATEDB NOCREATEROLE NOINHERIT;

-- Rol de auditor financiero
CREATE ROLE rol_auditor_fin
    NOLOGIN NOSUPERUSER NOCREATEDB NOCREATEROLE NOINHERIT;

-- Rol de backup (lectura completa para pg_dump)
CREATE ROLE rol_backup
    NOLOGIN NOSUPERUSER NOCREATEDB NOCREATEROLE NOINHERIT;
GRANT pg_read_all_data TO rol_backup;

-- Rol de solo lectura (monitoreo/BI)
CREATE ROLE rol_solo_lectura
    NOLOGIN NOSUPERUSER NOCREATEDB NOCREATEROLE NOINHERIT;
    \end{lstlisting}
\end{frame}

\begin{frame}{Código: Creación de Usuarios de Login}
    \begin{lstlisting}[title={Usuarios de Login con contraseña}]
CREATE ROLE artisync_app LOGIN PASSWORD 'App$ecure2026!'
    NOSUPERUSER INHERIT CONNECTION LIMIT 20 VALID UNTIL '2027-12-31';
GRANT rol_app_backend TO artisync_app;

CREATE ROLE artisync_admin LOGIN PASSWORD 'Adm!nStr0ng2026'
    NOSUPERUSER INHERIT CONNECTION LIMIT 3;
GRANT rol_administrador TO artisync_admin;

CREATE ROLE artisync_moderador LOGIN PASSWORD 'M0d3r@t0r2026'
    NOSUPERUSER INHERIT CONNECTION LIMIT 5;
GRANT rol_moderador TO artisync_moderador;

CREATE ROLE artisync_soporte LOGIN PASSWORD 'S0p0rt3$2026'
    NOSUPERUSER INHERIT CONNECTION LIMIT 5;
GRANT rol_soporte TO artisync_soporte;

CREATE ROLE artisync_auditor LOGIN PASSWORD 'Aud!t0r2026$'
    NOSUPERUSER INHERIT CONNECTION LIMIT 2;
GRANT rol_auditor_fin TO artisync_auditor;

CREATE ROLE artisync_backup LOGIN PASSWORD 'B@ckup$2026Secure'
    NOSUPERUSER INHERIT CONNECTION LIMIT 2;
GRANT rol_backup TO artisync_backup;

CREATE ROLE artisync_readonly LOGIN PASSWORD 'R3@d0nly2026!'
    NOSUPERUSER INHERIT CONNECTION LIMIT 10;
GRANT rol_solo_lectura TO artisync_readonly;
    \end{lstlisting}
\end{frame}

% ==============================================================================
% SECCIÓN 5: PRIVILEGIOS DEL SISTEMA
% ==============================================================================
\section{Privilegios del Sistema}

\begin{frame}{Privilegios del Sistema por Rol}
    \centering
    \scriptsize
    \begin{tabular}{>{\ttfamily\raggedright\arraybackslash}p{3.5cm} >{\raggedright\arraybackslash}p{3.8cm} >{\raggedright\arraybackslash}p{5.5cm}}
        \toprule
        \textbf{Rol} & \textbf{Privilegios} & \textbf{Justificación} \\
        \midrule
        rol\_administrador & \texttt{CREATEROLE}, \texttt{CREATEDB} & Crear roles personalizados y bases de prueba. \\
        \midrule
        rol\_app\_backend & \texttt{LOGIN}, \texttt{CONNECT} & Conexión JDBC; sin capacidad administrativa. \\
        \midrule
        rol\_moderador & \texttt{LOGIN}, \texttt{CONNECT} & Solo conexión; DML parcial. \\
        \midrule
        rol\_soporte & \texttt{LOGIN}, \texttt{CONNECT} & Solo conexión y consultas de lectura. \\
        \midrule
        rol\_auditor\_fin & \texttt{LOGIN}, \texttt{CONNECT} & Solo lectura sobre finanzas. \\
        \midrule
        rol\_backup & \texttt{LOGIN}, \texttt{pg\_read\_all\_data} & Lectura total para \texttt{pg\_dump}. \\
        \midrule
        rol\_solo\_lectura & \texttt{LOGIN}, \texttt{CONNECT} & Conexión mínima para BI. \\
        \bottomrule
    \end{tabular}

    \vspace{0.3cm}
    \begin{alertblock}{\faInfoCircle\ Nota}
        PostgreSQL no tiene \texttt{BACKUP DATABASE}. Se usa \texttt{pg\_read\_all\_data} como equivalente funcional.
    \end{alertblock}
\end{frame}

% ==============================================================================
% SECCIÓN 6: PRIVILEGIOS SOBRE OBJETOS
% ==============================================================================
\section{Privilegios sobre Objetos}

\begin{frame}{Privilegios: Backend y Administrador}
    \begin{columns}[T]
        \begin{column}{0.48\textwidth}
            \begin{block}{\texttt{rol\_app\_backend} --- DML Completo}
                \scriptsize
                \begin{tabular}{ll}
                    \toprule
                    \textbf{Privilegio} & \textbf{Objetos} \\
                    \midrule
                    \texttt{SELECT} & 37 tablas \\
                    \texttt{INSERT} & 37 tablas \\
                    \texttt{UPDATE} & 37 tablas \\
                    \texttt{DELETE} & 37 tablas \\
                    \texttt{USAGE} & Todas las secuencias \\
                    \texttt{EXECUTE} & Funciones PL/pgSQL \\
                    \texttt{USAGE} & \texttt{SCHEMA public} \\
                    \bottomrule
                \end{tabular}
            \end{block}
        \end{column}
        \begin{column}{0.48\textwidth}
            \begin{block}{\texttt{rol\_administrador} --- ALL + DDL}
                \scriptsize
                Hereda todo de \texttt{rol\_app\_backend} y añade:
                \begin{tabular}{ll}
                    \toprule
                    \textbf{Privilegio} & \textbf{Objetos} \\
                    \midrule
                    \texttt{ALL} & Todas las tablas \\
                    \texttt{CREATE} & \texttt{SCHEMA public} \\
                    \texttt{EXECUTE} & Todas las funciones \\
                    \texttt{TRIGGER} & Todas las tablas \\
                    \bottomrule
                \end{tabular}
            \end{block}
        \end{column}
    \end{columns}
\end{frame}

\begin{frame}{Privilegios: Moderador y Soporte}
    \begin{columns}[T]
        \begin{column}{0.48\textwidth}
            \begin{block}{\texttt{rol\_moderador}}
                \scriptsize
                \begin{tabular}{lp{3.5cm}}
                    \toprule
                    \textbf{Priv.} & \textbf{Tablas} \\
                    \midrule
                    \texttt{S} & \texttt{usuarios}, \texttt{perfiles\_creadores}, \texttt{portafolio\_items}, \texttt{servicios}, \texttt{likes\_portafolio} \\
                    \midrule
                    \texttt{S,U} & \texttt{certificados\_ia} \\
                    \midrule
                    \texttt{S,U,D} & \texttt{comentarios\_portafolio}, \texttt{sorteos} \\
                    \bottomrule
                \end{tabular}
            \end{block}
        \end{column}
        \begin{column}{0.48\textwidth}
            \begin{block}{\texttt{rol\_soporte} --- Solo Lectura}
                \scriptsize
                \texttt{SELECT} únicamente sobre:
                \begin{itemize}
                    \item \texttt{usuarios}, \texttt{usuario\_roles}
                    \item \texttt{sesiones\_usuario}
                    \item \texttt{pedidos}, \texttt{historial\_estados\_pedido}
                    \item \texttt{tickets\_revision}
                    \item \texttt{salas\_chat}, \texttt{mensajes}
                    \item \texttt{notificaciones\_sistema}
                    \item \texttt{servicios}, \texttt{contratos}
                \end{itemize}
            \end{block}
        \end{column}
    \end{columns}
\end{frame}

\begin{frame}{Privilegios: Auditor, Backup y Solo Lectura}
    \begin{columns}[T]
        \begin{column}{0.32\textwidth}
            \begin{block}{\scriptsize \texttt{rol\_auditor\_fin}}
                \scriptsize
                \texttt{SELECT} sobre:
                \begin{itemize}
                    \item \texttt{contratos}
                    \item \texttt{plantillas\_contrato}
                    \item \texttt{pagos\_garantia}
                    \item \texttt{transacciones\_pago}
                    \item \texttt{pedidos}
                    \item \texttt{entregables\_finales}
                \end{itemize}
            \end{block}
        \end{column}
        \begin{column}{0.32\textwidth}
            \begin{block}{\scriptsize \texttt{rol\_backup}}
                \scriptsize
                \begin{itemize}
                    \item \texttt{pg\_read\_all\_data}
                    \item Lectura completa de todo el esquema
                    \item Para ejecutar \texttt{pg\_dump}
                    \item Sin capacidad de modificación
                \end{itemize}
            \end{block}
        \end{column}
        \begin{column}{0.32\textwidth}
            \begin{block}{\scriptsize \texttt{rol\_solo\_lectura}}
                \scriptsize
                \begin{itemize}
                    \item \texttt{SELECT} sobre todas las tablas
                    \item \texttt{USAGE} en \texttt{SCHEMA public}
                    \item Para Grafana, Metabase y dashboards
                \end{itemize}
            \end{block}
        \end{column}
    \end{columns}

    \vspace{0.3cm}
    \centering
    \begin{exampleblock}{\faCheckCircle\ Principio Aplicado}
        Cada rol accede \textbf{exclusivamente} a las tablas requeridas por su función con los permisos \textbf{estrictamente necesarios}.
    \end{exampleblock}
\end{frame}

\begin{frame}{Matriz Resumen de Privilegios}
    \centering
    \scriptsize
    \begin{tabular}{l c c c c c c c}
        \toprule
        \textbf{Rol / Módulo} & \textbf{Seguridad} & \textbf{Perfiles} & \textbf{Catálogo} & \textbf{Pedidos} & \textbf{Finanzas} & \textbf{Comunic.} & \textbf{Social} \\
        \midrule
        \texttt{rol\_app\_backend}   & SIUD & SIUD & SIUD & SIUD & SIUD & SIUD & SIUD \\
        \texttt{rol\_administrador}  & SIUD+DDL & SIUD+DDL & SIUD+DDL & SIUD+DDL & SIUD+DDL & SIUD+DDL & SIUD+DDL \\
        \texttt{rol\_moderador}      & S & S,U & S & --- & --- & --- & S,U,D \\
        \texttt{rol\_soporte}        & S & S & --- & S & --- & S & --- \\
        \texttt{rol\_auditor\_fin}   & --- & --- & --- & S & S & --- & --- \\
        \texttt{rol\_backup}         & S* & S* & S* & S* & S* & S* & S* \\
        \texttt{rol\_solo\_lectura}  & S & S & S & S & S & S & S \\
        \bottomrule
    \end{tabular}

    \vspace{0.2cm}
    \tiny S=SELECT, I=INSERT, U=UPDATE, D=DELETE, DDL=CREATE/ALTER/DROP, S*=pg\_read\_all\_data
\end{frame}

% ==============================================================================
% SECCIÓN 7: CÓDIGO DE PRIVILEGIOS
% ==============================================================================
\section{Código de Privilegios}

\begin{frame}{Código: Privilegios del Backend}
    \begin{lstlisting}[title={DML completo para rol\_app\_backend}]
-- Revocar privilegios por defecto
REVOKE ALL ON SCHEMA public FROM PUBLIC;
REVOKE ALL ON ALL TABLES IN SCHEMA public FROM PUBLIC;

-- USAGE del esquema
GRANT USAGE ON SCHEMA public TO rol_app_backend;
GRANT CREATE ON SCHEMA public TO rol_administrador;

-- SELECT, INSERT, UPDATE, DELETE sobre TODAS las tablas
GRANT SELECT, INSERT, UPDATE, DELETE
    ON ALL TABLES IN SCHEMA public TO rol_app_backend;

-- USAGE sobre todas las secuencias (SERIAL)
GRANT USAGE, SELECT
    ON ALL SEQUENCES IN SCHEMA public TO rol_app_backend;

-- EXECUTE sobre funciones PL/pgSQL
GRANT EXECUTE
    ON ALL FUNCTIONS IN SCHEMA public TO rol_app_backend;

-- Privilegios por defecto para tablas FUTURAS
ALTER DEFAULT PRIVILEGES IN SCHEMA public
    GRANT SELECT, INSERT, UPDATE, DELETE ON TABLES TO rol_app_backend;
ALTER DEFAULT PRIVILEGES IN SCHEMA public
    GRANT USAGE, SELECT ON SEQUENCES TO rol_app_backend;
ALTER DEFAULT PRIVILEGES IN SCHEMA public
    GRANT EXECUTE ON FUNCTIONS TO rol_app_backend;
    \end{lstlisting}
\end{frame}

\begin{frame}{Código: Privilegios Moderador, Soporte y Auditor}
    \begin{lstlisting}[title={Privilegios granulares por rol}]
-- === ROL MODERADOR ===
GRANT SELECT ON usuarios, perfiles_creadores, portafolio_items,
    portafolios, servicios, likes_portafolio,
    estados_verificacion, habilidades TO rol_moderador;
GRANT SELECT, UPDATE ON certificados_ia TO rol_moderador;
GRANT SELECT, UPDATE, DELETE ON comentarios_portafolio TO rol_moderador;
GRANT SELECT, UPDATE, DELETE ON sorteos TO rol_moderador;

-- === ROL SOPORTE (solo lectura operativa) ===
GRANT SELECT ON usuarios, usuario_roles, roles, perfiles_creadores,
    sesiones_usuario, pedidos, historial_estados_pedido,
    tickets_revision, salas_chat, mensajes, documentos_adjuntos,
    notificaciones_sistema, servicios, contratos TO rol_soporte;

-- === ROL AUDITOR FINANCIERO (solo lectura financiera) ===
GRANT SELECT ON contratos, plantillas_contrato, pagos_garantia,
    transacciones_pago, pedidos, entregables_finales,
    servicios, usuarios TO rol_auditor_fin;

-- === ROL SOLO LECTURA (SELECT global) ===
GRANT SELECT ON ALL TABLES IN SCHEMA public TO rol_solo_lectura;
ALTER DEFAULT PRIVILEGES IN SCHEMA public
    GRANT SELECT ON TABLES TO rol_solo_lectura;
    \end{lstlisting}
\end{frame}

% ==============================================================================
% SECCIÓN 8: HARDENING
% ==============================================================================
\section{Hardening y Verificación}

\begin{frame}{Hardening: Revocación de Privilegios Peligrosos}
    \begin{lstlisting}[title={Revocación explícita y endurecimiento}]
-- Revocar CONNECT a la BD para PUBLIC
REVOKE CONNECT ON DATABASE artisyncbd FROM PUBLIC;

-- Otorgar CONNECT solo a roles autorizados
GRANT CONNECT ON DATABASE artisyncbd TO rol_app_backend;
GRANT CONNECT ON DATABASE artisyncbd TO rol_administrador;
GRANT CONNECT ON DATABASE artisyncbd TO rol_moderador;
GRANT CONNECT ON DATABASE artisyncbd TO rol_soporte;
GRANT CONNECT ON DATABASE artisyncbd TO rol_auditor_fin;
GRANT CONNECT ON DATABASE artisyncbd TO rol_backup;
GRANT CONNECT ON DATABASE artisyncbd TO rol_solo_lectura;

-- Revocar CREATE sobre esquema para roles no admin
REVOKE CREATE ON SCHEMA public FROM rol_app_backend;
REVOKE CREATE ON SCHEMA public FROM rol_moderador;
REVOKE CREATE ON SCHEMA public FROM rol_soporte;

-- Revocar TRUNCATE para evitar eliminaciones masivas
REVOKE TRUNCATE ON ALL TABLES IN SCHEMA public
    FROM rol_app_backend;
REVOKE TRUNCATE ON ALL TABLES IN SCHEMA public
    FROM rol_moderador;
    \end{lstlisting}
\end{frame}

\begin{frame}{Consultas de Verificación}
    \begin{lstlisting}[title={Validar la configuracion de seguridad}]
-- Listar todos los roles creados
SELECT rolname, rolsuper, rolcreaterole, rolcreatedb,
       rolcanlogin, rolconnlimit
FROM pg_roles
WHERE rolname LIKE 'artisync_%' OR rolname LIKE 'rol_%';

-- Verificar membresias de roles
SELECT r.rolname AS usuario, m.rolname AS miembro_de
FROM pg_auth_members am
JOIN pg_roles r ON am.member = r.oid
JOIN pg_roles m ON am.roleid = m.oid
WHERE r.rolname LIKE 'artisync_%';

-- Verificar privilegios sobre tablas
SELECT grantee, table_name, privilege_type
FROM information_schema.table_privileges
WHERE table_schema = 'public' AND grantee LIKE 'rol_%'
ORDER BY grantee, table_name, privilege_type;

-- Verificar que PUBLIC no tiene acceso
SELECT grantee, table_name, privilege_type
FROM information_schema.table_privileges
WHERE table_schema = 'public' AND grantee = 'PUBLIC';
    \end{lstlisting}
\end{frame}

% ==============================================================================
% SECCIÓN 9: DOBLE CAPA
% ==============================================================================
\section{Seguridad en Doble Capa}

\begin{frame}{Relación: Roles de BD vs Roles de Aplicación}
    \begin{columns}[T]
        \begin{column}{0.48\textwidth}
            \begin{block}{\faDatabase\ Capa de Base de Datos}
                \scriptsize
                \begin{itemize}
                    \item Roles PostgreSQL (\texttt{rol\_*})
                    \item Controlan acceso a nivel de \textbf{tablas y esquema}
                    \item Privilegios: \texttt{SELECT}, \texttt{INSERT}, etc.
                    \item Defensa perimetral de datos
                \end{itemize}
            \end{block}
        \end{column}
        \begin{column}{0.48\textwidth}
            \begin{block}{\faCode\ Capa de Aplicación}
                \scriptsize
                \begin{itemize}
                    \item Roles en tabla \texttt{roles} (ADMIN, CREADOR, CLIENTE, MODERADOR, SOPORTE, AUDITOR\_FINANCIERO)
                    \item Spring Security + JWT
                    \item Permisos granulares: \texttt{USUARIO\_VER}, \texttt{ROL\_GESTIONAR}, etc.
                    \item Filtrado por lógica de negocio
                \end{itemize}
            \end{block}
        \end{column}
    \end{columns}

    \vspace{0.3cm}
    \centering
    \scriptsize
    \begin{tabular}{lll}
        \toprule
        \textbf{Rol Aplicación} & \textbf{Rol BD} & \textbf{Relación} \\
        \midrule
        ADMIN & \texttt{rol\_administrador} & Acceso directo + vía API \\
        CREADOR & \texttt{rol\_app\_backend} & Solo vía API (mediado) \\
        CLIENTE & \texttt{rol\_app\_backend} & Solo vía API (mediado) \\
        MODERADOR & \texttt{rol\_moderador} & Panel directo + API \\
        SOPORTE & \texttt{rol\_soporte} & Herramienta de soporte \\
        AUDITOR\_FINANCIERO & \texttt{rol\_auditor\_fin} & Reportes directos \\
        \bottomrule
    \end{tabular}
\end{frame}

% ==============================================================================
% CONCLUSIONES
% ==============================================================================
\section{Conclusiones}

\begin{frame}{Conclusiones}
    \begin{enumerate}
        \item[\faUsers] Se identificaron \textbf{9 tipos de usuarios}: 6 de negocio y 3 técnicos.

        \vspace{0.2cm}
        \item[\faLayerGroup] Se definieron \textbf{7 roles de grupo} en PostgreSQL con jerarquía clara de herencia.

        \vspace{0.2cm}
        \item[\faCog] Los \textbf{privilegios de sistema} se asignaron restrictivamente: solo el administrador posee \texttt{CREATEROLE} y \texttt{CREATEDB}.

        \vspace{0.2cm}
        \item[\faLock] Los \textbf{privilegios sobre objetos} siguen el principio de mínimo privilegio: cada rol accede exclusivamente a sus tablas necesarias.

        \vspace{0.2cm}
        \item[\faShieldAlt] Se implementó \textbf{seguridad en doble capa}: base de datos (PostgreSQL) + aplicación (Spring Security con JWT).
    \end{enumerate}

    \vspace{0.3cm}
    \begin{exampleblock}{\faCheckCircle\ Resultado}
        \centering Script idempotente, documentado y verificable listo para despliegue en desarrollo, staging y producción.
    \end{exampleblock}
\end{frame}

% --- CIERRE ---
\begin{frame}
    \centering
    \vspace{1.5cm}
    {\Huge\bfseries\textcolor{uteqgreen}{¡Gracias!}}

    \vspace{0.8cm}
    {\Large Preguntas y Respuestas}

    \vspace{1cm}
    \small
    Bone Niurca $\cdot$ Carvajal Johan $\cdot$ Figueroa Bryan $\cdot$ Rios Jhon

    \vspace{0.3cm}
    \scriptsize
    Universidad Técnica Estatal de Quevedo\\
    Facultad de Ciencias de la Ingeniería

    \vspace{0.5cm}
    \faGithub\ ArtiSync --- \faDatabase\ PostgreSQL 16
\end{frame}

\end{document}
